\ifdefined\iscombined
	\lecturetitle{Bisimilarity}
	\setcounter{section}{5}
\else
\documentclass{article}
\usepackage{etoolbox}
\usepackage{enumitem}
\usepackage{amsthm}
\usepackage{comment}
\usepackage{amsmath}
\usepackage{amsfonts}
\usepackage{sectsty}
\usepackage{hyperref}
\usepackage{fancyhdr}
\usepackage[top=1in,bottom=1in,left=1in,right=1in]{geometry}
\usepackage{float}
\usepackage{graphicx}
\usepackage{tabularx}
\usepackage{mathtools}

\title{Lecture 5 - Bisimilarity}
\author{Matthew Farah, \href{mailto:farahm11@mcmaster.ca}{$\langle$farahm11@mcmaster.ca$\rangle$}, 400468251}
\date{\today}

\hypersetup{
	colorlinks=true,
	linkcolor=blue,
	filecolor=magenta,
	urlcolor=blue,
	citecolor=blue,
	anchorcolor=blue
}

\fancyhead{}
\fancyhead[L]{Matthew Farah, \href{mailto:farahm11@mcmaster.ca}{$\langle$farahm11@mcmaster.ca$\rangle$}, 400468251}
\fancyhead[R]{Lecture 5 - Bisimilarity}

\newcounter{example}
\newcounter{partctr}[example]

\newenvironment{example}{
	\refstepcounter{example}
	\setcounter{partctr}{0}
	\vspace{1em}
	\large\noindent\textbf{Ex.}\vspace{0.5em}
}{\vspace{1.5em}}

\newenvironment{parts}{
	\begin{enumerate}[
		label=\textbf{(\alph*)},
		leftmargin=3.3em,
		itemsep=0.4em,
		before=\large
	]
	\setcounter{enumi}{\value{partctr}}
}{
	\setcounter{partctr}{\value{enumi}}
	\end{enumerate}
}

\newenvironment{solution}{
	\vspace{0.3cm}\par\noindent\large\textbf{Solution:}\par\vspace{0.3em}
}{\vspace{0.8em}}

\renewcommand{\thesubsection}{\thesection.\alph{subsection}}
\renewcommand{\thesubsubsection}{\thesection.\alph{subsection}.\Roman{subsubsection}}
\makeatletter

\sectionfont{\fontsize{12}{12}\bfseries}
\subsectionfont{\fontsize{10}{12}\normalfont}
\subsubsectionfont{\fontsize{10}{12}\normalfont}

\newcommand{\parts@seriesname}{parts@\theexample}

\begin{document}

\maketitle
\pagestyle{fancy}

\fi

\begin{itemize}
	\item The \emph{traces} of a process \texttt{P} to some state \texttt{p} in its LTS representation is the set of all sequences of actions which can be taken to reach \texttt{p}.
	\begin{itemize}
		\item The traces of a process to its starting state is denoted by Trace(\texttt{P}).
	\end{itemize}
	\item Two processes, \texttt{P}, \texttt{Q} are said to be \emph{equivalent} if for some third process \texttt{R}, $(\mathtt{R} \; || \; \mathtt{P}) = (\mathtt{R} \; || \; \mathtt{Q})$.
	\begin{itemize}
		\item \texttt{P} and \texttt{Q} are denoted as being equivalent by $\mathtt{P} \equiv \mathtt{Q}$.
	\end{itemize}
	\item If \texttt{P} and \texttt{Q} are sequential or deterministic processes, then $\mathtt{P} \equiv \mathtt{Q} \iff \text{Traces}(\mathtt{P}) = \text{Traces}(\mathtt{Q})$.
	\item \underline{However}, in any other case, it is not sufficient for the traces of the two processes to be equal, since under parallel composition, two processes with the same traces can behave differently.
	\begin{itemize}
		\item For example, for non-deterministic processes \texttt{P} and \texttt{Q} where $\text{Trace}(\mathtt{P}) = \text{Trace}(\mathtt{Q})$, it can be true that $(P \; || \; R)$ can deadlock whereas $(Q \; || \; R)$ does not for some process \texttt{R}.
	\end{itemize}
	\item Hence, in the context of concurrent systems, equivalence is redefined as \emph{bisimilarity}.
	\item Consider the LTS representations of \texttt{P} and \texttt{Q}:
	\begin{itemize}
		\item Two states $\mathtt{p} \in \mathtt{P}$ and $\mathtt{q} \in \mathtt{Q}$ are said to be bisimilar if every action at \texttt{p} can be executed at \texttt{q}, and likewise if every action at \texttt{q} can be performed at \texttt{p}.
		\begin{itemize}
			\item \texttt{p} and \texttt{q} being bisimilar is denoted $\mathtt{p} \approx \mathtt{q}$.
			\item Bisimilarity is both associative and commutative.
			\item Important to note that a single state in \texttt{P} can be bisimilar to multiple states in \texttt{Q} and vice versa.
		\end{itemize}
		\item Two processes \texttt{P} and \texttt{Q} are said to be bisimilar if, for every state $\mathtt{p} \in \mathtt{P}$ reachable by traces \emph{t}, those some traces reach a state $\mathtt{q} \in \mathtt{Q}$ which is bisimilar to \texttt{p}.
	\end{itemize}
	\item \texttt{P} and \texttt{Q} are denoted as being bisimilar by $\mathtt{P} \approx \mathtt{Q}$.
	\begin{itemize}
		\item $\mathtt{P} \approx \mathtt{Q} \implies (\mathtt{P} \; || \; \mathtt{R}) = (\mathtt{P} \; || \; \mathtt{R})$ for all processes \texttt{R}.
	\end{itemize}
	%TODO: Provide example
	\item Bisimilarity cannot distinguish between two processes which can exhbit simultaneity or interleaving
\end{itemize}

%TODO: Justify solutions

\newpage

\begin{example}
	Determine whether the following LTSs are bisimilar.
\end{example}

\begin{parts}
	\item
\end{parts}

\begin{figure}[H]
	\centering
	\includegraphics[width=\linewidth/2]{~/Desktop/University/Courses/SFWRENG-3BB4-Coursework/Lecture-Notes/Lecture-5/imgs/Example-1-LTSs.png}
	\caption{LTSs for part (a).}
\end{figure}

\begin{solution}
	They are bisimilar
\end{solution}

\newpage

\begin{parts}
	\item
\end{parts}

\begin{figure}[H]
	\centering
	\includegraphics[width=\linewidth]{~/Desktop/University/Courses/SFWRENG-3BB4-Coursework/Lecture-Notes/Lecture-5/imgs/Example-2-LTSs.png}
	\caption{LTSs for part (b).}
\end{figure}

\begin{solution}
	They are not bisimilar
\end{solution}

\ifdefined\iscombined
\newpage
\else
  \end{document}
\fi
